% -----------------------------------------------------------
% Brochure for Ciao 1.14 announcement (for ICLP'06 and later)
% -----------------------------------------------------------
\documentclass[letterpaper]{article}

\usepackage{twocolumn}
\usepackage[letterpaper,inner=1in,outer=1in,top=1in,bottom=1in]{geometry}
\usepackage{ciaocopy}
\usepackage{graphicx} 
%\usepackage{fullpage}
% \textwidth 15.8cm
%\textwidth 16.3cm
%\tolerance=10000

\usepackage{cmbright}
% \usepackage{charter}
% \usepackage{times}
% \usepackage{helvet}
% \usepackage[scaled]{helvet}
% \usepackage{utopia}
% \usepackage{avant}
% \fontfamily{phv}
% \selectfont

\makeatletter
\def\subsection{\@startsection{subsection}{2}{\z@}{-1.0ex plus -1ex minus 
    -.2ex}{0.7ex plus .2ex}{\normalsize\bf}}
\makeatother

\renewcommand{\baselinestretch}{1.0}
\newcommand{\myitems}{
  \renewcommand{\baselinestretch}{0.5}
  \parsep 0pt \parindent 0pt \itemsep -4pt \topsep -3pt }

%%
% psfigTeX macros: combined TeX/LaTeX source
%
% Previous copyright below applies.  University of Maryland changes are
% Copyright (c) 1989 University of Maryland
% Department of Computer Science.  All rights reserved.
% Permission to copy these changes for any purpose is hereby granted
% so long as this copyright notice remains intact.
%
% psfig.tex and psfig.sty are derived automatically from this file,
% by including the appropriate one of either the TeX or LaTeX specific
% lines (marked with `%tex:' and `%latex').
%
% -------------------------------------------
%
% All software, documentation, and related files in this distribution of
% psfig/tex are Copyright (c) 1987 Trevor J. Darrell
%
% Permission is granted for use and non-profit distribution of psfig/tex
% providing that this notice be clearly maintained, but the right to
% distribute any portion of psfig/tex for profit or as part of any commercial
% product is specifically reserved for the author.
%
% Psfig/tex version 1.1
%
% $Header: /usr/src/local/tex/local/mctex/psfig/RCS/psfig.src,v 3.1 89/08/30 03:42:01 chris Exp $
% based on: psfig.tex,v 1.8 87/07/25 13:21:09 trevor Exp
%

\catcode`\@=11\relax
\newwrite\@unused
\def\typeout#1{{\let\protect\string\immediate\write\@unused{#1}}}

\typeout{psfig: version 1.1 (MC-TeX)}

\newdimen\ps@dim

%\def\psglobal#1{\ps@typeout
%  {psfig: including #1 globally}\immediate\special{ps:plotfile #1 global}}
%\def\psfiginit{\ps@typeout{psfiginit}\psglobal{/usr/lib/ps/figtex.pro}}

% @psdo control structure -- similar to Latex @for.
% I redefined these with different names so that psfig can
% be used with TeX as well as LaTeX, and so that it will not
% be vunerable to future changes in LaTeX's internal
% control structure.
%
\def\@nnil{\@nil}
\def\@empty{}
\def\@psdonoop#1\@@#2#3{}
\def\@psdo#1:=#2\do#3{\edef\@psdotmp{#2}\ifx\@psdotmp\@empty \else
\expandafter\@psdoloop#2,\@nil,\@nil\@@#1{#3}\fi}
\def\@psdoloop#1,#2,#3\@@#4#5{\def#4{#1}\ifx #4\@nnil \else
#5\def#4{#2}\ifx #4\@nnil \else#5\@ipsdoloop #3\@@#4{#5}\fi\fi}
\def\@ipsdoloop#1,#2\@@#3#4{\def#3{#1}\ifx #3\@nnil
\let\@nextwhile=\@psdonoop \else
#4\relax\let\@nextwhile=\@ipsdoloop\fi\@nextwhile#2\@@#3{#4}}

%
%

\def\psdraft{\def\@psdraft{0}}
\def\psfull{\def\@psdraft{100}}
\psfull

\def\ps@eat#1{}
\def\pssilent{\let\ps@typeout\ps@eat}
\def\psverbose{\let\ps@typeout\typeout}
\psverbose

\newif\if@prologfile
\newif\if@postlogfile
%%% These are for the option list.
%%% A specification of the form a = b maps to calling \@p@@sa{b}
\newif\if@bbllx
\newif\if@bblly
\newif\if@bburx
\newif\if@bbury
\newif\if@height
\newif\if@width
\newif\if@rheight
\newif\if@rwidth
\newif\if@clip
\def\@p@@sclip#1{\@cliptrue}
\def\@p@@sfile#1{\def\@p@sfile{#1}}
\def\@p@@sfigure#1{\def\@p@sfile{#1}}
\def\@p@@sbbllx#1{%\ps@typeout{bbllx is #1}
	\@bbllxtrue
	\ps@dim=#1
	\edef\@p@sbbllx{\number\ps@dim}
}
\def\@p@@sbblly#1{
	%\ps@typeout{bblly is #1}
	\@bbllytrue
	\ps@dim=#1
	\edef\@p@sbblly{\number\ps@dim}
}
\def\@p@@sbburx#1{
	%\ps@typeout{bburx is #1}
	\@bburxtrue
	\ps@dim=#1
	\edef\@p@sbburx{\number\ps@dim}
}
\def\@p@@sbbury#1{
	%\ps@typeout{bbury is #1}
	\@bburytrue
	\ps@dim=#1
	\edef\@p@sbbury{\number\ps@dim}
}
\def\@p@@sheight#1{
	\@heighttrue
	\ps@dim=#1
	\edef\@p@sheight{\number\ps@dim}
	%\ps@typeout{Height is \@p@sheight}
}
\def\@p@@swidth#1{
	%\ps@typeout{Width is #1}
	\@widthtrue
	\ps@dim=#1
	\edef\@p@swidth{\number\ps@dim}
}
\def\@p@@srheight#1{
	%\ps@typeout{Reserved height is #1}
	\@rheighttrue
	\ps@dim=#1
	\edef\@p@srheight{\number\ps@dim}
}
\def\@p@@srwidth#1{
	%\ps@typeout{Reserved width is #1}
	\@rwidthtrue
	\ps@dim=#1
	\edef\@p@srwidth{\number\ps@dim}
}
\def\@p@@sprolog#1{\@prologfiletrue\def\@prologfileval{#1}}
\def\@p@@spostlog#1{\@postlogfiletrue\def\@postlogfileval{#1}}
\def\@p@@ssilent#1{\pssilent}

\def\@cs@name#1{\csname #1\endcsname}
\def\@setparms#1=#2,{\@cs@name{@p@@s#1}{#2}}
%
% initialize the defaults (size the size of the figure)
%
\def\ps@init@parms{
	\@bbllxfalse \@bbllyfalse
	\@bburxfalse \@bburyfalse
	\@heightfalse \@widthfalse
	\@rheightfalse \@rwidthfalse
	\def\@p@sbbllx{}\def\@p@sbblly{}
	\def\@p@sbburx{}\def\@p@sbbury{}
	\def\@p@sheight{}\def\@p@swidth{}
	\def\@p@srheight{}\def\@p@srwidth{}
	\def\@p@sfile{}
	\def\@p@scost{10}
	\def\@sc{}
	\@prologfilefalse
	\@postlogfilefalse
	\@clipfalse
}
%
% Go through the options setting things up.
%
\def\parse@ps@parms#1{\@psdo\@psfiga:=#1\do{\expandafter\@setparms\@psfiga,}}

%
% Compute bb height and width
%
\newif\ifno@bb
\newif\ifnot@eof
\newread\ps@stream
\def\bb@missing{
	\ps@typeout{psfig: searching \@p@sfile \space  for bounding box}
	\openin\ps@stream=\@p@sfile
	\no@bbtrue
	\not@eoftrue
	\catcode`\%=12
	\loop
		\read\ps@stream to \line@in
		\global\toks200=\expandafter{\line@in}
		\ifeof\ps@stream \not@eoffalse \fi
		%\ps@typeout{looking at :: \the\toks200}
		\@bbtest{\toks200}
		\if@bbmatch\not@eoffalse\expandafter\bb@cull\the\toks200\fi
	\ifnot@eof \repeat
	\catcode`\%=14
}
\catcode`\%=12
\newif\if@bbmatch
\def\@bbtest#1{\expandafter\@a@\the#1%%BoundingBox:\@bbtest\@a@}
\long\def\@a@#1%%BoundingBox:#2#3\@a@{\ifx\@bbtest#2\@bbmatchfalse\else\@bbmatchtrue\fi}
\long\def\bb@cull#1 #2 #3 #4 #5 {
	\ps@dim=#2 bp\edef\@p@sbbllx{\number\ps@dim}
	\ps@dim=#3 bp\edef\@p@sbblly{\number\ps@dim}
	\ps@dim=#4 bp\edef\@p@sbburx{\number\ps@dim}
	\ps@dim=#5 bp\edef\@p@sbbury{\number\ps@dim}
	\no@bbfalse
}
\catcode`\%=14
%
\def\compute@bb{
	\no@bbfalse
	\if@bbllx \else \no@bbtrue \fi
	\if@bblly \else \no@bbtrue \fi
	\if@bburx \else \no@bbtrue \fi
	\if@bbury \else \no@bbtrue \fi
	\ifno@bb \bb@missing \fi
	\ifno@bb \typeout{FATAL ERROR: no bb supplied or found}
		\no-bb-error
	\fi
	%
	\count203=\@p@sbburx
	\count204=\@p@sbbury
	\advance\count203 by -\@p@sbbllx
	\advance\count204 by -\@p@sbblly
	\edef\@bbw{\number\count203}
	\edef\@bbh{\number\count204}
	%\ps@typeout{bbh = \@bbh, bbw = \@bbw}
}
%
% \in@hundreds performs #1 * (#2 / #3) correct to the hundreds,
%	then leaves the result in @result
%
\def\in@hundreds#1#2#3{\count240=#2 \count241=#3
	\count100=\count240	% 100 is first digit #2/#3
	\divide\count100 by \count241
	\count101=\count100
	\multiply\count101 by \count241
	\advance\count240 by -\count101
	\multiply\count240 by 10
	\count101=\count240	%101 is second digit of #2/#3
	\divide\count101 by \count241
	\count102=\count101
	\multiply\count102 by \count241
	\advance\count240 by -\count102
	\multiply\count240 by 10
	\count102=\count240	% 102 is the third digit
	\divide\count102 by \count241
	\count200=#1\count205=0
	\count201=\count200
		\multiply\count201 by \count100
	 	\advance\count205 by \count201
	\count201=\count200
		\divide\count201 by 10
		\multiply\count201 by \count101
		\advance\count205 by \count201
		%
	\count201=\count200
		\divide\count201 by 100
		\multiply\count201 by \count102
		\advance\count205 by \count201
		%
	\edef\@result{\number\count205}
}
\def\compute@wfromh{
	% computing : width = height * (bbw / bbh)
	\in@hundreds{\@p@sheight}{\@bbw}{\@bbh}
	%\ps@typeout{ \@p@sheight * \@bbw / \@bbh, = \@result }
	\edef\@p@swidth{\@result}
	%\ps@typeout{w from h: width is \@p@swidth}
}
\def\compute@hfromw{
	% computing : height = width * (bbh / bbw)
	\in@hundreds{\@p@swidth}{\@bbh}{\@bbw}
	%\ps@typeout{ \@p@swidth * \@bbh / \@bbw = \@result }
	\edef\@p@sheight{\@result}
	%\ps@typeout{h from w: height is \@p@sheight}
}
\def\compute@handw{
	\if@height \if@width \else \compute@wfromh \fi
	\else \if@width \compute@hfromw \else
		\edef\@p@sheight{\@bbh}
		\edef\@p@swidth{\@bbw} \fi
	\fi
}
\def\compute@resv{
	\if@rheight \else \edef\@p@srheight{\@p@sheight} \fi
	\if@rwidth \else \edef\@p@srwidth{\@p@swidth} \fi
}
%
% Compute any missing values
\def\compute@sizes{
	\compute@bb
	\compute@handw
	\compute@resv
}
%
% \psfig
% usage : \psfig{file=, height=, width=, bbllx=, bblly=, bburx=, bbury=,
%			rheight=, rwidth=, clip=}
%
% "clip=" is a switch and takes no value, but the `=' must be preset.
\def\psfig#1{\vbox{
    \ps@init@parms
    \parse@ps@parms{#1}
    \compute@sizes
    %
    \ifnum\@p@scost<\@psdraft
	\ps@typeout{psfig: including \@p@sfile}
	%
	\special{ps::[begin] \@p@swidth \space \@p@sheight \space
		\@p@sbbllx \space \@p@sbblly \space
		\@p@sbburx \space \@p@sbbury \space
		startTexFig \space}
	\if@clip
		\ps@typeout{(clip)}
		\special{ps:: \@p@sbbllx \space \@p@sbblly \space
			\@p@sbburx \space \@p@sbbury \space
			doclip \space}
	\fi
	\if@prologfile \special{ps: plotfile \@prologfileval \space} \fi
	\special{ps: plotfile \@p@sfile \space}
	\if@postlogfile \special{ps: plotfile \@postlogfileval \space} \fi
	\special{ps::[end] endTexFig \space}
	% Create the vbox to reserve the space for the figure
	\vbox to\@p@srheight true sp{\hbox to\@p@srwidth true sp{\hfil}\vfil}
    \else
	% draft figure, just reserve the space and print the
	% path name.
	\vbox to\@p@srheight true sp{\vss
	    \hbox to\@p@srwidth true sp{\hss\@p@sfile\hss}\vss}
    \fi
}}

\catcode`\@=12\relax


%\input{/home/clip/Papers/clip_description/logo.tex}
%\newcommand{\ciao}{\psfig{figure=ciao_s.ps,width=2.5cm}}
\newcommand{\ciao}{%
    \includegraphics[viewport=205 372 362 420,scale=0.75]{ciao_s}%
    }

\begin{document}
\textheight = 670pt
\pagestyle{empty}

\twocolumn[
\noindent
\centerline{
\hbox{
  \begin{Huge}
    {\bf The 
    \hspace*{-0.09\textwidth}
      \begin{minipage}[c]{0.3\textwidth}
        \ciao
      \end{minipage}
    \hspace*{-0.06\textwidth}
      1.14 System}\\
  \end{Huge}
}
}

\centerline{Ciao is \emph{free software} distributed under GNU
  licenses.}  
\vspace{0.5\baselineskip}
]

\noindent
Ciao is a modern, multiparadigm programming language with an advanced
programming environment. It includes a full high-performance
\emph{ISO-standard} Prolog system, as well as very significant
extensions: its novel modular design allows both \emph{restricting}
and \emph{extending} the language through libraries. This makes it
possible to work with \emph{fully declarative subsets} of Prolog and
also to \emph{extend} these subsets (or ISO-Prolog) both syntactically
and semantically.  Most importantly, these restrictions and extensions
can be activated separately on each program module so that several
extensions can coexist in the same application for different modules.

Ciao currently includes libraries for functional logic programming,
higher-order (with predicate abstractions), constraints,
object-oriented programming, feature terms (records), persistence,
several control rules, concurrency, distributed execution (agents),
concurrent built-in database, and parallel execution.  Libraries also
support WWW programming, sockets, external interfaces, etc.\ (see the
\textbf{Main Features} section.)

\emph{Programming in the large} is supported thanks to a robust
module/class system, module-based automatic incremental compilation
(with no need for makefiles), an assertion language for declaring
(\emph{optional}) program properties (including types, modes,
determinacy, non-failure, cost, ...), automatic static inference and
static/dynamic checking of such assertions, program optimization and
parallelization, and powerful automatic documentation generation (the
latter tasks performed by the \emph{CiaoPP preprocessor} and the
\emph{LPdoc autodocumenter}).  It is the novel modular design of Ciao
which enables modular program development, effective global program
analysis, and static debugging and optimization via source to source
program transformation.

\emph{Programming in the small} is supported by having reduced size
executables, which only include those builtins and libraries used by
the program, and by supporting Prolog \emph{scripts}.

The compiler allows generating very easily several forms of
architecture-independent and stand-alone executables. Program modules
can be
%% compiled into compact bytecode and 
linked statically, dynamically, or autoloaded. The executables
generated are very competitive in both performance and size with all
current commercial and academic Prolog systems.  
%% Optimizing compilation to 100\% native code is undergoing work.

The programming environment also offers a rich graphical interface
(with direct access to top-level/debugger, preprocessor, and
autodocumenter, based on Emacs), an embeddable source-level debugger
with breakpoints, and execution visualization tools.

\medskip
% \hrule

\subsection*{Why do we call it Ciao?}
% After reading this brochure 
The sharp reader may have already seen the logic behind the 'Ciao
Prolog' name. Ciao means both \emph{hello} and \emph{goodbye}, and
Ciao Prolog intends to be on one hand a really excellent, freely
available ISO-Prolog system, for both academic and industrial use
(including introducing users to Prolog and to classical constraint and
logic programming) --- the \emph{hello} part.  But Ciao is also a
new-generation, multiparadigm programming language and program
development system which goes well beyond Prolog and other classical
logic programming languages. And it has the advantage (when compared
to other systems) that it does so while allowing full Prolog
compatibility when needed.

%% Ciao is a \emph{next generation} logic programming environment.

\medskip\hrule

\subsection*{Main Features of Ciao:}
\begin{itemize}
  \myitems
  % General:
\item Efficient, high-performance, bytecode-based engine, with garbage
  collection, unbound precision integer arithmetic, built-in
  concurrency capabilities, and many other features.
  % TA:
\item  % Full 
  Extensive compliance with 
  % International Standard ISO/IEC 13211-1
  ISO-Prolog.
\item Generation of \textbf{multi-architecture} executables: Linux,
  {\Large $*$}NIX, Mac OS X, Win32.
\item User-friendly installation on all platforms.
\item Advanced, integrated development environment with
  integrated source-level debugging, syntax highlighting, 
  on-line help, etc. % (\texttt{emacs}-based).
  %% \item Source code debugger (with code highlighting).
\item Debugger embeddable in executables.
  % TA:
\item Exception handling.
\item Source code autodocumenter (with a menu-based interface).
  % 
  % Module system/Libraries:
\item New generation, robust module system.
\item Modular clp(R) / clp(Q).
\item Higher-order syntax and predicate abstractions.
  % \item Library of higher-order list predicates.
  % \item Libraries for and-fair breadth-first and iterative deepening.
\item Extensive, built in, and modular code expansion facilities
  (macros) with operators local to modules.
\item Attributed variables, DCGs.
\item Backwards compatibility libraries (DEC-10 IO, Quintus-like
  internal database, etc.). 
\item Libraries of (commented) types, modes, and other properties to
  be used in assertions (for debugging and documentation generation).
\item Assertion-based declaration of meta-predicates.
  % 
\item Flexible customization of library paths and path aliases.
  % Programming styles:
\item Several execution strategies available: Andorra, breadth-first,
  iterative deepening.
\item Object-oriented extensions.
\item Bidirectional foreign interfaces: C, Java, TclTk, ProVRML.
\item Interface to SQL, relational databases.
\item Concurrency / multithreading. % engine primitives (\&Prolog-like). 
  % TA:
\item Full thread support in Linux / Unix / Mac OS X / Win32.
\item Access to operating system resources.
  % TA:
\item Platform independent socket communication.
  % TA:
\item Platform independent Prolog Makefiles.
\item Delay predicates (when/2, freeze/1).
\item Active modules (distr./agent programming).
\item Remote loading of modules.
\item Web/Internet programming: \emph{PiLLoW} library.

\item Programmer-transparent 
  persistent predicates (files or relational database storage).
\item Execution of fuzzy logic-based programs.
  % 
  % Engine/Compiler features:
\item Ciao CGI executables under IIS.
\item Number of clauses/predicates essentially unbound.
\item Unbound atom size.
\item Fast creation of new unique atoms.
\item Fast writing/reading (marshalling and unmarshalling) of terms.
\item Compressed object code/executables. % on demand. 
\item Fast compilation and startup. 
\item Incremental stand-alone compiler with separate compilation.
\item Automatic (re)compilation of foreign files.
\item Attributed variables.
\item Extensive, up to date documentation in multiple formats.
\end{itemize}
\vspace*{-3mm}
... and many more!
\vspace*{1mm}

\hrule

% \subsection*{What's new / improved in 1.14?}
\subsection*{Upcoming in Ciao 1.14:}

{\small\emph{Final beta 1.13.X versions are being posted in
    preparation for 1.14.0 release, including:}}

\begin{itemize}
  \myitems
\item New libraries:
  \begin{itemize}
    \myitems\vspace*{-0.75em}
  \item POSIX / shell regular expression search and replacement.
  \item Lazy evaluation (composable with functional syntax).
  \item Execution profiler.
  \item Backtrackable global variables.
  \item Package to help writing metainterpreters.
  \end{itemize}
\item Improved libraries: 
  \begin{itemize}
    \myitems\vspace*{-0.75em}
  \item Association lists with self-adapting data structure.
  \item Interface to O.S.\ facilities enriched.
  \item New publishing possibilities for active modules.
  \item Functional notation split into two different packages.
%%     with automatic evaluation of arithmetic functors
%%     (\texttt{functional}) or without it (\texttt{fsyntax}).
  \item General improvements to the MySQL interface.
  \item Fuzzy Prolog enhanced.
  \item ``\texttt{file\_utils}'' library enriched.
  \item \texttt{make} library enhanced.
  \item \texttt{statistics} provides more information.
  \item \texttt{clpq} and \texttt{clpr} can dump and
    project constraints.
  \end{itemize}
%% Companion tools
\item The \texttt{LPdoc} automatic documenter is now included in the
  default Ciao packaging.
\item Tighter integration with CiaoPP, the Ciao program analysis and
  transformation tool.
%% User experience
\item Emacs mode enhanced.
\item Improved compatibility with XEmacs.
%%
\item New configuration and installation scheme: \texttt{Makefile}s
  replaced by Ciao \texttt{Makefile.pl}s.
\item % Many 
      Bug fixes both in the abstract machine and in the Prolog libraries.
\item Infinite terms printed in equational form.
% \item Floating point numbers accurately displayed.
\item Added exceptions to more libraries.
\item Preliminary ports to Mac OS X / i386, XSCale / Linux, Power PC /
  Linux.
\item More modules annotated with % type \& mode 
      assertions.
\item Many existing assertions improved/corrected.
\item More properties added to the standard set.
\item Assertion library improved.
%% Interfaces
\item C code can  be embedded in Ciao files.
\item Added DLL support in Windows.
%% Extensions and support
\item Added CHR implementation (thanks to T.\ Schrijvers).
\item Support for multiattributes.
%% Utilities
\item XDR schema to HTML forms utility.
\item Interface to GnuPlot.
\item Interface to \texttt{dot} (graph representation library).
\item Bidirectional list traversal library.
\item More flexible, richer stream handling.


%% \item Interface to the Parma Polyhedra Library, if available.
%% \item Tighter integration with the Ciao program analysis and
%%   transformation tool, CiaoPP.
%\item Compile-time macro facilities.

  %% \hrule

  %% % \item Double-click startup of programming environment.
  %%   menus.
  %% \item \texttt{XEmacs} compatibility improved (thanks to A. Rigo).
  %% \item File aliasing to internal streams added.
  %% \item The \texttt{make} library has been improved.
  %% \item PiLLoW library improved in many ways.
  %% \item Tcl/Tk library improved and made more robust.
  %% \item Improved assertions package.
  %% \item Atom-to-term conversion library improved.
  %% \item Exceptions in active modules improved.
  %% \item DaVinci error processing improved.
  %% \item Persistent predicates improved.
  %% \item File locking capabilities included.
  %% \item New input/output facilities added to socket-related library.
  %% \item Added treatment of\- operators and 
  %%   \texttt{module:pred} calls to pretty-printer.
  %% \item Updated reporting of read syntax errors.
  %% \item Unbound length atoms now supported in all cases.
  %% \item C interface \texttt{.h} files reachable through a more standard
  %%   location (thanks to R. Bagnara).
  %% \item Compatibility with command line options of newer \texttt{gcc}
  %%   versions.
  %% \item Solved conflicts in predicate reexportation.
  %% \item Replication of clauses in some cases corrected (thanks to
  %%   S. Craig).
  %% \item Got rid of several SEGV problems.
  %% \item The number of significant decimal digits to be printed is now
  %%   computed accurately.
  %% \item The reader has improved its behavior w.r.t. non-terminated
  %%   comments.
  %%   
  %% \item The \texttt{xref} library treats void references now.
  %% \item Foreign predicates are now automatically declared as
  %%   implementation-defined ones.
  %% \item New utilities to build module frontiers.
  %% \item Facts defined in external files can now be automatically cached
  %%   on-demand.
  %% \item And, in general, new features in many libraries.
\end{itemize}

\hrule

\subsection*{Additional contributions:\\ (in development or beta state)}
\begin{itemize}
  \myitems
\item Ongoing integration with Emacs 22-beta.
\item Distributed Hash Tables library.
\item Port to Ultra Sparc / Solaris.
\item Improved CLP(FD).
%% \item XML querying and transformation to Prolog.
%% 
\item Library implementing Lambda Prolog (thanks to J.\ Chanco \& G.\
  Nadathur).  
\end{itemize}
%% \noindent
%% \textbf{Ciao Prolog 1.10 will be available in short} 

%% \noindent
%% Release candidate versions are being posted at the Ciao Prolog WWW
%% site.  We very much appreciate feedback on Ciao Prolog!

%% \medskip
%% \hrule

\hrule

\subsection*{Contact / download info:}
\texttt{http://cliplab.org}\\ 
\texttt{http://ciaoprolog.org}\\ 
\texttt{ciao@clip.dia.fi.upm.es}\\ [2mm]
The CLIP Group\\
Technical University of Madrid, Spain\\
University of New Mexico, USA\\
%% Facultad de Inform\'{a}tica -- UPM\\
%% E-28660 Boadilla del Monte, Madrid, SPAIN

\hrule

\end{document}
